% GENERATED FILE. Edit canonical lessons, then run npm run generate:books.
% canonical-source-hash: fnv1a64:373842a9117c0add
% canonical-lessons: MW-C169-povno, MW-C169-belno, MW-C169-ukalno, MW-C169-talno, MW-C169-chakhno

\chapter{Making Rotis}
\label{ch:mw-making-rotis}
\hlchaptermodality{169}

\begin{chapteropening}
\textbf{By the end of this chapter:} \emph{I can say how the dough is kneaded and the food is cooked.}

Complete the last lesson of chapter 169: I can say how the dough is kneaded and the food is cooked.
\end{chapteropening}

\section[\texorpdfstring{\mw{पोवणो} (povṇo)}{पोवणो (povṇo)}]{\mw{पोवणो} (povṇo) — to knead (dough); to make rotis by hand}

\label{lesson:MW-C169-povno}



\begin{quote}
Before the new one: say the Marwadi for holy basil, then the Marwadi for to water plants.
\end{quote}

\subsection*{You'll want to know: \mw{पोवणो}}
\textbf{\mw{पोवणो}} — \emph{povṇo} — \textquotedblleft{}to knead (dough); to make rotis by hand\textquotedblright{}.

\begin{quote}
\mw{मां} \mw{दिनूगै} \mw{रोटी} \mw{पोवै} \mw{है।}
\end{quote}

\begin{quote}
\emph{mā̃ dinūgai roṭī povai hai.}
\end{quote}

\begin{quote}
Mother makes the rotis by hand early in the morning.
\end{quote}

Past (I …): \mw{म्हैं} \mw{रोटी} \mw{पोई}, \emph{mhaĩ roṭī poī}.

\begin{cousinweb}
\textbf{\mw{पोवणो}} (\emph{povṇo}) is inherited, by CDIAL's account, from Sanskrit \emph{pratudati}, \textquotedblleft{}cuts through, pierces\textquotedblright{}; CDIAL lists Old Marwari \emph{povaï} twice, as \textquotedblleft{}kneads (bread)\textquotedblright{} and as \textquotedblleft{}threads (beads)\textquotedblright{} (CDIAL 8627).
\end{cousinweb}

\begin{grammarlens}[title={What you've built}]
One more word for saying what you do.
\end{grammarlens}

\subsection*{Your turn}
\begin{itemize}
\raggedright
  \item \emph{Say it:} \emph{povṇo}
  \item \emph{Say it:} \emph{povṇo}, once more
  \item \emph{Say it:} \emph{mā̃ dinūgai roṭī povai hai.} — Mother makes the rotis by hand early in the morning.
  \item \emph{Recall:} say the Marwadi for holy basil, then the Marwadi for to water plants, then say \emph{povṇo} again
\end{itemize}

\begin{tcolorbox}[breakable,colback=teal!4,colframe=teal!35!black,title={Before you move on}]
What does \mw{पोवणो} mean? (\textquotedblleft{}To knead (dough); to make rotis by hand\textquotedblright{}.) Say it once more.
\end{tcolorbox}
\section[\texorpdfstring{\mw{बेलणो} (belṇo)}{बेलणो (belṇo)}]{\mw{बेलणो} (belṇo) — to roll out (dough)}

\label{lesson:MW-C169-belno}



\begin{quote}
Before the new one: say the Marwadi for to water plants, then the Marwadi for to knead.
\end{quote}

\subsection*{You'll want to know: \mw{बेलणो}}
\textbf{\mw{बेलणो}} — \emph{belṇo} — \textquotedblleft{}to roll out (dough)\textquotedblright{}.

\begin{quote}
\mw{हूं} \mw{रोज} \mw{रोटी} \mw{बेलूं} \mw{हूं।}
\end{quote}

\begin{quote}
\emph{hū̃ roj roṭī belū̃ hū̃.}
\end{quote}

\begin{quote}
I roll out the rotis every day.
\end{quote}

Past (I …): \mw{म्हैं} \mw{रोटी} \mw{बेली}, \emph{mhaĩ roṭī belī}.

\begin{cousinweb}
\textbf{\mw{बेलणो}} (\emph{belṇo}) is inherited from Sanskrit \emph{vellayati}, \textquotedblleft{}causes to shake, kneads\textquotedblright{}; CDIAL lists Hindi \emph{belnā}, \textquotedblleft{}to roll out\textquotedblright{} (CDIAL 12121). CDIAL gives no Marwari form in this entry. The rolling pin, \textbf{\mw{बेलण}}, is its tool.
\end{cousinweb}

\begin{grammarlens}[title={What you've built}]
One more word for saying what you do.
\end{grammarlens}

\subsection*{Your turn}
\begin{itemize}
\raggedright
  \item \emph{Say it:} \emph{belṇo}
  \item \emph{Say it:} \emph{belṇo}, once more
  \item \emph{Say it:} \emph{hū̃ roj roṭī belū̃ hū̃.} — I roll out the rotis every day.
  \item \emph{Recall:} say the Marwadi for to water plants, then the Marwadi for to knead, then say \emph{belṇo} again
\end{itemize}

\begin{tcolorbox}[breakable,colback=teal!4,colframe=teal!35!black,title={Before you move on}]
What does \mw{बेलणो} mean? (\textquotedblleft{}To roll out (dough)\textquotedblright{}.) Say it once more.
\end{tcolorbox}
\section[\texorpdfstring{\mw{उकळणो} (ukaḷṇo)}{उकळणो (ukaḷṇo)}]{\mw{उकळणो} (ukaḷṇo) — to boil}

\label{lesson:MW-C169-ukalno}



\begin{quote}
Before the new one: say the Marwadi for to knead, then the Marwadi for to roll out.
\end{quote}

\subsection*{You'll want to know: \mw{उकळणो}}
\textbf{\mw{उकळणो}} — \emph{ukaḷṇo} — \textquotedblleft{}to boil\textquotedblright{}.

\begin{quote}
\mw{हांडी} \mw{मांय} \mw{पाणी} \mw{उकळै} \mw{है।}
\end{quote}

\begin{quote}
\emph{hā̃ḍī mā̃y pāṇī ukḷai hai.}
\end{quote}

\begin{quote}
The water boils in the pot.
\end{quote}

Past (it …): \mw{वो} \mw{दूध} \mw{उकळ्यो}, \emph{vo dūdh ukaḷyo}.

\begin{cousinweb}
\textbf{\mw{उकळणो}} (\emph{ukaḷṇo}) comes from a reconstructed \emph{*utkalati}, \textquotedblleft{}goes out or up\textquotedblright{}; CDIAL lists Hindi \emph{ukalnā}, \textquotedblleft{}to boil over\textquotedblright{}, and Gujarati \emph{ukaḷvũ} (CDIAL 1716). CDIAL gives no Marwari form in this entry.
\end{cousinweb}

\begin{grammarlens}[title={What you've built}]
One more word for saying what you do.
\end{grammarlens}

\subsection*{Your turn}
\begin{itemize}
\raggedright
  \item \emph{Say it:} \emph{ukaḷṇo}
  \item \emph{Say it:} \emph{ukaḷṇo}, once more
  \item \emph{Say it:} \emph{hā̃ḍī mā̃y pāṇī ukḷai hai.} — The water boils in the pot.
  \item \emph{Recall:} say the Marwadi for to knead, then the Marwadi for to roll out, then say \emph{ukaḷṇo} again
\end{itemize}

\begin{tcolorbox}[breakable,colback=teal!4,colframe=teal!35!black,title={Before you move on}]
What does \mw{उकळणो} mean? (\textquotedblleft{}To boil\textquotedblright{}.) Say it once more.
\end{tcolorbox}
\section[\texorpdfstring{\mw{तळणो} (taḷṇo)}{तळणो (taḷṇo)}]{\mw{तळणो} (taḷṇo) — to fry}

\label{lesson:MW-C169-talno}



\begin{quote}
Before the new one: say the Marwadi for to roll out, then the Marwadi for to boil.
\end{quote}

\subsection*{You'll want to know: \mw{तळणो}}
\textbf{\mw{तळणो}} — \emph{taḷṇo} — \textquotedblleft{}to fry\textquotedblright{}.

\begin{quote}
\mw{मां} \mw{कड़ाही} \mw{मांय} \mw{पापड़} \mw{तळै} \mw{है।}
\end{quote}

\begin{quote}
\emph{mā̃ kaṛāhī mā̃y pāpaṛ taḷai hai.}
\end{quote}

\begin{quote}
Mother fries papads in the wok.
\end{quote}

Past (I …): \mw{म्हैं} \mw{पापड़} \mw{तळ्यो}, \emph{mhaĩ pāpaṛ taḷyo}.

\begin{cousinweb}
\textbf{\mw{तळणो}} (\emph{taḷṇo}) comes from a reconstructed \emph{*talati}, \textquotedblleft{}fries\textquotedblright{}; CDIAL lists Hindi \emph{talnā} and Gujarati \emph{taḷvũ} (CDIAL 5736). CDIAL gives no Marwari form in this entry.
\end{cousinweb}

\begin{grammarlens}[title={What you've built}]
One more word for saying what you do.
\end{grammarlens}

\subsection*{Your turn}
\begin{itemize}
\raggedright
  \item \emph{Say it:} \emph{taḷṇo}
  \item \emph{Say it:} \emph{taḷṇo}, once more
  \item \emph{Say it:} \emph{mā̃ kaṛāhī mā̃y pāpaṛ taḷai hai.} — Mother fries papads in the wok.
  \item \emph{Recall:} say the Marwadi for to roll out, then the Marwadi for to boil, then say \emph{taḷṇo} again
\end{itemize}

\begin{tcolorbox}[breakable,colback=teal!4,colframe=teal!35!black,title={Before you move on}]
What does \mw{तळणो} mean? (\textquotedblleft{}To fry\textquotedblright{}.) Say it once more.
\end{tcolorbox}
\section[\texorpdfstring{\mw{चाखणो} (chākhṇo)}{चाखणो (chākhṇo)}]{\mw{चाखणो} (chākhṇo) — to taste}

\label{lesson:MW-C169-chakhno}



\begin{quote}
Before the new one: say the Marwadi for to boil, then the Marwadi for to fry.
\end{quote}

\subsection*{You'll want to know: \mw{चाखणो}}
\textbf{\mw{चाखणो}} — \emph{chākhṇo} — \textquotedblleft{}to taste\textquotedblright{}.

\begin{quote}
\mw{हूं} \mw{दाल} \mw{चाखूं} \mw{हूं।}
\end{quote}

\begin{quote}
\emph{hū̃ dāl chākhū̃ hū̃.}
\end{quote}

\begin{quote}
I taste the dal.
\end{quote}

Past (I …): \mw{म्हैं} \mw{दाल} \mw{चाखी}, \emph{mhaĩ dāl chākhī}.

\begin{cousinweb}
\textbf{\mw{चाखणो}} (\emph{chākhṇo}) comes from a reconstructed \emph{*cakṣati}, \textquotedblleft{}tastes\textquotedblright{}; CDIAL lists Hindi \emph{cākhnā} and Gujarati \emph{cākhvũ} (CDIAL 4557). CDIAL gives no Marwari form in this entry.
\end{cousinweb}

\begin{grammarlens}[title={What you've built}]
One more word for saying what you do.
\end{grammarlens}

\subsection*{Your turn}
\begin{itemize}
\raggedright
  \item \emph{Say it:} \emph{chākhṇo}
  \item \emph{Say it:} \emph{chākhṇo}, once more
  \item \emph{Say it:} \emph{hū̃ dāl chākhū̃ hū̃.} — I taste the dal.
  \item \emph{Recall:} say the Marwadi for to boil, then the Marwadi for to fry, then say \emph{chākhṇo} again
\end{itemize}

\begin{tcolorbox}[breakable,colback=teal!4,colframe=teal!35!black,title={Before you move on}]
What does \mw{चाखणो} mean? (\textquotedblleft{}To taste\textquotedblright{}.) Say it once more.
\end{tcolorbox}
